% Del diagrama de flujo a los bloques de MakeCode — programar es traducir.
%
% Diagrama propio en TikZ (sin librerías extra). Se tiñe con el color de la línea
% (milcGradeDark = naranja de Robótica). Muestra, lado a lado, cómo cada figura
% del diagrama de flujo del módulo 2 se convierte en un bloque de MakeCode: el
% óvalo → «por siempre», el rectángulo → una acción, el rombo → «si … si no».
% Deja ver que el módulo 3 no inventa: traduce el plano del módulo 2.
%
% Se incrusta vía \guiaDiagrama desde el bloque `recursos:` del YAML.
\begin{tikzpicture}[scale=1,font=\sffamily,>=stealth]

  % ── Encabezados de las dos columnas ─────────────────────────────────
  \node[font=\scriptsize\bfseries,milcNegro] at (1.6,6.2) {diagrama de flujo (módulo 2)};
  \node[font=\scriptsize\bfseries,milcNegro] at (8.4,6.2) {bloques de MakeCode (módulo 3)};
  \draw[milcNegro!30,line width=.4pt] (5,6.5) -- (5,-0.3);

  % ── Fila 1: óvalo → «por siempre» ────────────────────────────────────
  % óvalo
  \fill[milcGradeDark] (0.5,5.2) rectangle (2.7,5.7);
  \fill[milcGradeDark] (0.5,5.45) circle (0.25);
  \fill[milcGradeDark] (2.7,5.45) circle (0.25);
  \node[white,font=\scriptsize\bfseries] at (1.6,5.45) {inicio (repetir)};
  % bloque «por siempre»
  \fill[milcTurquesa!85] (6.4,5.05) rectangle (10.4,5.85);
  \node[white,font=\scriptsize\bfseries,anchor=west] at (6.6,5.45) {por siempre \{ \ldots \}};
  \draw[->,milcNegro!60,line width=.7pt] (3.3,5.45) -- (6.3,5.45);

  % ── Fila 2: rectángulo → variable = lectura ─────────────────────────
  \fill[milcGradeDark!85] (0.5,3.9) rectangle (2.7,4.6);
  \node[white,font=\scriptsize\bfseries,align=center] at (1.6,4.25) {leer luz\\(sensor)};
  \fill[milcTurquesa!85] (6.4,3.85) rectangle (10.6,4.65);
  \node[white,font=\scriptsize,anchor=west] at (6.6,4.25) {\ttfamily luz = nivel de luz};
  \draw[->,milcNegro!60,line width=.7pt] (3.3,4.25) -- (6.3,4.25);

  % ── Fila 3: rombo → «si … si no» ─────────────────────────────────────
  \fill[milcMostaza] (1.6,3.5) -- (2.9,2.6) -- (1.6,1.7) -- (0.3,2.6) -- cycle;
  \draw[milcNegro!55,line width=.4pt] (1.6,3.5) -- (2.9,2.6) -- (1.6,1.7) -- (0.3,2.6) -- cycle;
  \node[milcNegro,font=\scriptsize\bfseries,align=center] at (1.6,2.6) {¿luz $<$\\umbral?};
  \fill[milcTurquesa!85] (6.4,1.55) rectangle (10.8,3.65);
  \node[white,font=\scriptsize,anchor=west,align=left] at (6.6,2.6)
    {\ttfamily si (luz < umbral) \{\\ \ \ mostrar luna\\ \} si no \{\\ \ \ mostrar sol\\ \}};
  \draw[->,milcNegro!60,line width=.7pt] (3.3,2.6) -- (6.3,2.6);

  % ── Fila 4: actuadores ───────────────────────────────────────────────
  \fill[milcGradeDark!85] (0.2,0.7) rectangle (1.5,1.4);
  \node[white,font=\scriptsize\bfseries,align=center] at (0.85,1.05) {mostrar\\luna};
  \fill[milcGradeDark!85] (1.8,0.7) rectangle (3.1,1.4);
  \node[white,font=\scriptsize\bfseries,align=center] at (2.45,1.05) {mostrar\\sol};
  \node[font=\scriptsize\itshape,milcNegro!70,anchor=west,text width=4.1cm] at (6.4,1.0)
    {cada acción del diagrama es un bloque «mostrar»; el actuador es la matriz de LED};

  % (El mensaje «traducir, no inventar» va en el pie de figura, para no repetirlo.)

\end{tikzpicture}
